\chapter{Structural Credit Models}\label{ch:rc:structural-credit-models}

On 4 May 2005 General Motors' share price rose 18\% after Kirk Kerkorian offered to buy
a large stake in the company: his company Tracinda then tendered for up to 28 million shares at 31
dollars each. The next day Standard \& Poor's cut the company's debt two
notches, to BB, below investment grade, and its bond and default-swap spreads widened.
Trades that bet on the prices of the company's debt and equity staying in line with each
other lost on both legs, on consecutive days. The model behind many such trades is the
subject of this chapter: a firm defaults when the value of its assets falls short of its
debt, so that equity is an option on the firm and the prices of its shares and of its
credit are two views of the same asset value. The \omterm{def:rc:reduced-form-credit:reduced}{reduced-form model} of chapter~13
takes the hazard rate as given; the \omterm{def:rc:structural-credit-models:structural}{structural model} explains it, and links it to the
equity market, which is what makes it useful and what made May 2005 expensive.

\section{Equity as a call on the firm}

\begin{definition}[Structural model]\label{def:rc:structural-credit-models:structural}
A \emph{structural model}\index{structural model} of credit models the value $V_t$ of a
firm's assets and makes default an event defined by $V$ and the firm's liabilities:
assets below the debt at its maturity, or assets touching a barrier. Default is then
predictable from the asset value, and the prices of equity and debt follow from it.
\end{definition}

\begin{definition}[Merton model]\label{def:rc:structural-credit-models:merton}
The \emph{Merton model}\index{Merton model} takes assets following a geometric Brownian
motion with volatility $\sigma_V$ and one zero-coupon debt of face $D$ due at $T$. At $T$
the shareholders receive $\max(V_T-D,0)$ and the creditors $\min(V_T,D)$. Equity is a
European call on the assets struck at the debt, priced by the Black--Scholes formula,
and debt is the assets minus the equity:
\[
E_0 = V_0N(d_1)-De^{-rT}N(d_2),\qquad B_0 = V_0-E_0,\qquad
d_{1,2} = \frac{\ln(V_0/D)+(r\pm\tfrac12\sigma_V^2)T}{\sigma_V\sqrt T}.
\]
\end{definition}

Equivalently, debt is a riskless bond minus a put on the assets struck at $D$: the
creditors have sold the shareholders the right to hand over the firm instead of repaying.
The credit spread of the debt is $s = -\ln(B_0/D)/T - r$, and the risk-neutral
\omterm{def:rc:reduced-form-credit:pd}{probability of default} is $N(-d_2)$. \Cref{fig:rc:structural-credit-models:payoff} shows
the two claims as functions of the asset value.

\begin{omfigure}
\begin{tikzpicture}
\begin{axis}[width=11.5cm,height=5.4cm,
  xlabel={asset value $V_0$ (USD billion)},ylabel={value (USD billion)},
  tick label style={font=\small},label style={font=\small},
  xmin=0,xmax=80,ymin=0,ymax=45,grid=major,grid style={gray!25},
  legend columns=3,legend style={font=\footnotesize,at={(0.5,-0.25)},anchor=north,draw=none,fill=none,/tikz/every even column/.append style={column sep=8pt}},
  table/col sep=comma]
\addplot[omThm,very thick] table[x=v,y=equity]{figdata/rates-credit-risk/14-structural-credit-models/equity_debt.csv};
\addplot[omDef,very thick] table[x=v,y=debt]{figdata/rates-credit-risk/14-structural-credit-models/equity_debt.csv};
\addplot[omExo,thick,dashed] table[x=v,y=payoff]{figdata/rates-credit-risk/14-structural-credit-models/equity_debt.csv};
\legend{equity today,debt today,equity at maturity}
\end{axis}
\end{tikzpicture}

\omcaption{Equity and debt of the chapter's firm (debt face USD~40 billion due in five
years, asset volatility 15.6\%) as functions of the asset value. Equity is a call on the
assets; debt rises towards the riskless value of the face and falls one-for-one with the
assets when they are low. Data: the chapter's tutorial.}\label{fig:rc:structural-credit-models:payoff}
\end{omfigure}

The asset value and its volatility are not observed. Equity is: its market value $E_0$
and its volatility $\sigma_E$, implied from options or estimated from returns. By Itô's
formula on $E = C(V)$, $\sigma_EE_0 = N(d_1)\sigma_VV_0$, and with the pricing equation this
gives two equations for the two unknowns.

\omcode{code/firm/structural/firm_structural.py}{65}{81}{Backing out asset value and asset volatility from equity value and equity volatility (Newton in two unknowns).}\label{lst:rc:structural-credit-models:invert}

\begin{example}[A leveraged firm]\label{ex:rc:structural-credit-models:firm}
An illustrative firm has equity worth USD~10 billion with a volatility of 50\%, and debt
of face USD~40 billion treated as due in five years; $r = 4\%$. The inversion gives assets
of USD~40.69 billion with a volatility of 15.6\%: the equity is more than three times as
volatile as the assets because it is levered. The debt is worth USD~30.69 billion, a yield spread
of 130 basis points, and the risk-neutral probability that the assets end below the face
is 32.7\%.
\end{example}

\section{Distance to default}

\begin{definition}[Distance to default]\label{def:rc:structural-credit-models:dd}
The \emph{distance to default}\index{distance to default} of a firm over a horizon $T$ is
the number of standard deviations by which the expected log asset value exceeds the
default point:
\[
\mathrm{DD} = \frac{\ln(V_0/D)+(\mu-\tfrac12\sigma_V^2)T}{\sigma_V\sqrt T},
\]
with $\mu$ the real-world asset drift. In the \omterm{def:rc:structural-credit-models:merton}{Merton model} the real-world \omterm{def:rc:reduced-form-credit:pd}{probability of
default} is $N(-\mathrm{DD})$; in practice the distance is mapped to default frequencies
through a history of defaults, because asset returns are not normal.
\end{definition}

\begin{example}[Two probabilities]\label{ex:rc:structural-credit-models:dd}
With an asset drift of 8\%, the firm's five-year \omterm{def:rc:structural-credit-models:dd}{distance to default} is 1.02 and the
real-world \omterm{def:rc:reduced-form-credit:pd}{probability of default} 15.4\%, against the risk-neutral 32.7\% of
\cref{ex:rc:structural-credit-models:firm}. The gap is the \omterm{def:rc:reduced-form-credit:premium}{default risk premium} of
chapter~13 in structural form: under the risk-neutral measure the assets drift at $r$
instead of $\mu$.
\end{example}

The \omterm{def:rc:structural-credit-models:merton}{Merton model}'s term structure of spreads depends on the quasi-leverage
$d = De^{-rT}/V_0$ (\cref{fig:rc:structural-credit-models:merton}). A firm with low leverage
has spreads that start at zero and rise with maturity, since it takes time for the assets
to fall to the debt; a firm near the default point has a humped curve; an insolvent one
($d>1$) has spreads falling with maturity, since only time can rescue it.

\begin{omfigure}
\begin{tikzpicture}
\begin{axis}[width=11.5cm,height=5.4cm,
  xlabel={maturity (years)},ylabel={Merton spread (bp)},
  tick label style={font=\small},label style={font=\small},
  xmin=0,xmax=10,ymin=0,ymax=2000,grid=major,grid style={gray!25},
  restrict y to domain=0:2100,
  legend columns=3,legend style={font=\footnotesize,at={(0.5,-0.25)},anchor=north,draw=none,fill=none,/tikz/every even column/.append style={column sep=8pt}},
  table/col sep=comma]
\addplot[omMeth,very thick] table[x=t,y=d05]{figdata/rates-credit-risk/14-structural-credit-models/merton_ts.csv};
\addplot[omDef,very thick,dashed] table[x=t,y=d09]{figdata/rates-credit-risk/14-structural-credit-models/merton_ts.csv};
\addplot[omThm,very thick,dotted] table[x=t,y=d12]{figdata/rates-credit-risk/14-structural-credit-models/merton_ts.csv};
\legend{$d = 0.5$,$d = 0.9$,$d = 1.2$}
\end{axis}
\end{tikzpicture}

\omcaption{Merton yield spreads against maturity at three quasi-leverages
$d = De^{-rT}/V_0$, asset volatility 25\%: rising for a safe firm, humped near the default
point, falling for a firm whose assets are worth less than its discounted debt. Data: the
chapter's tutorial.}\label{fig:rc:structural-credit-models:merton}
\end{omfigure}

\begin{remark}[Spreads that are too low]\label{rem:rc:structural-credit-models:low}
Merton spreads are low for safe firms and short maturities: a diffusion cannot reach the
default point in a few months, so short spreads are near zero while markets quote tens of
basis points. Remedies add jumps to the assets, uncertainty about the default point, or
a barrier that can be touched at any time, the next section.
\end{remark}

\section{First-passage models}

\begin{definition}[First-passage model, default barrier]\label{def:rc:structural-credit-models:fp}
A \emph{first-passage model}\index{first-passage model} defaults the firm the first time
its assets touch a \emph{default barrier}\index{default barrier} $B_t$ below the current
value, whenever that happens, instead of only at the debt's maturity. A barrier stands for
covenants that let creditors take over, or for the level at which the firm can no longer
refinance. The CreditGrades model (2002), published by RiskMetrics with three dealers, puts the barrier
at the average recovery on debt times the debt per share and makes that recovery lognormally uncertain,
so that default can come as a surprise and short spreads are not zero.
\end{definition}

\begin{proposition}[Survival to a constant barrier]\label{prop:rc:structural-credit-models:fp}
With risk-neutral asset dynamics $dV = rV\,dt+\sigma_VV\,dW$, a constant barrier $B<V_0$,
$x = \ln(V_0/B)$ and $\nu = r-\tfrac12\sigma_V^2$, the probability of no default up to $t$ is
\[
Q(t) = N\!\left(\frac{x+\nu t}{\sigma_V\sqrt t}\right)
- e^{-2\nu x/\sigma_V^2}\,N\!\left(\frac{-x+\nu t}{\sigma_V\sqrt t}\right).
\]
\end{proposition}
\begin{proof}
$\ln V_t-\ln V_0$ is a Brownian motion with drift $\nu$ and volatility $\sigma_V$; the
probability that it stays above $-x$ up to $t$ follows from the reflection principle and a
change of measure removing the drift (One Quant Book~4, chapter~2).
\end{proof}

\omcode{code/firm/structural/firm_structural.py}{93}{100}{The first-passage survival probability, in closed form.}\label{lst:rc:structural-credit-models:fp}

The survival function becomes a hazard curve on the premium dates, and chapter~13's legs
turn it into a model default-swap spread for each maturity, with a recovery set on the
default swap's convention (40\%) rather than read from the asset value at the barrier.

\begin{example}[The firm under a barrier]\label{ex:rc:structural-credit-models:fp}
With the barrier at 70\% of the debt face, the firm's model default-swap spreads are 64,
173, 216, 229 and 197 basis points at one, two, three, five and ten years
(\cref{fig:rc:structural-credit-models:fp}); its five-year survival probability is 82.3\%. A
market quote of 500 basis points at five years needs a barrier at 78.7\% of the face.
\end{example}

\begin{omfigure}
\begin{tikzpicture}
\begin{axis}[width=11.5cm,height=5.4cm,
  xlabel={maturity (years)},ylabel={model CDS spread (bp)},
  tick label style={font=\small},label style={font=\small},
  xmin=0,xmax=10,ymin=0,ymax=650,grid=major,grid style={gray!25},
  legend columns=3,legend style={font=\footnotesize,at={(0.5,-0.25)},anchor=north,draw=none,fill=none,/tikz/every even column/.append style={column sep=8pt}},
  table/col sep=comma]
\addplot[omMeth,very thick,mark=*,mark size=1.2pt] table[x=t,y=l60]{figdata/rates-credit-risk/14-structural-credit-models/fp_ts.csv};
\addplot[omDef,very thick,mark=square*,mark size=1.2pt] table[x=t,y=l70]{figdata/rates-credit-risk/14-structural-credit-models/fp_ts.csv};
\addplot[omThm,very thick,mark=triangle*,mark size=1.6pt] table[x=t,y=limp]{figdata/rates-credit-risk/14-structural-credit-models/fp_ts.csv};
\legend{barrier 60\% of face,barrier 70\%,barrier 78.7\% (fits 500 bp)}
\end{axis}
\end{tikzpicture}

\omcaption{Model default-swap spreads of the chapter's firm in the \omterm{def:rc:structural-credit-models:fp}{first-passage model},
for three barriers. Raising the barrier lifts the curve and moves its hump to shorter
maturities; the highest barrier fits a five-year quote of 500 basis points. Data: the
chapter's tutorial.}\label{fig:rc:structural-credit-models:fp}
\end{omfigure}

\begin{dated}{2026-09}{May 2005}\label{dat:rc:structural-credit-models:may2005}
The Bank for International Settlements' review of June 2005 reported that on 4 May 2005
General Motors' stock rose 18\% after an offer from Kirk Kerkorian to buy a large stake,
that on 5 May Standard \& Poor's cut GM by two notches to BB and Ford by one notch to
BB+, and that relative-value trades between GM's convertible bonds and a replicating
portfolio lost on both legs as the bond and default-swap spreads rose.
\end{dated}

\section{Capital-structure signals and arbitrage}

A \omterm{def:rc:structural-credit-models:structural}{structural model} maps an equity price to a credit spread. Traders use the map in both
directions: equity moves predict spread moves, and spreads far from the model's value
are candidate trades.

\begin{definition}[Capital-structure arbitrage]\label{def:rc:structural-credit-models:csa}
\emph{Capital-structure arbitrage}\index{capital-structure arbitrage} trades one
instrument of a firm's capital structure against another (default swaps or bonds
against equity or equity options) when their prices disagree with a \omterm{def:rc:structural-credit-models:structural}{structural model},
hedging with the model's sensitivity of one to the other: sell protection and short
equity when the market spread is above the model's, buy protection and buy equity when
it is below.
\end{definition}
Yu (2006) ran such a convergence strategy with CreditGrades on 135\,759 daily default-swap spreads of 261
North American names: single trades could lose heavily because spreads and equity prices were weakly
correlated, while an equally weighted portfolio of all trades earned Sharpe ratios similar to other
fixed-income arbitrage strategies.


\begin{method}[The equity hedge of a default swap]\label{met:rc:structural-credit-models:hedge}
Hold the asset volatility fixed. Move the equity value by a small amount, find the asset
value that reprices it, and recompute the model spread: this gives $\partial s/\partial E$.
Multiply by the position's \omterm{def:rc:reduced-form-credit:cs01}{CS01} (per unit of spread) to get its value change per unit of
equity value, and trade that much equity in the opposite direction.
\end{method}

\begin{example}[The hedge ratio]\label{ex:rc:structural-credit-models:hedge}
With the barrier at 70\% of the face, the model spread falls by 0.47 basis points for
each USD~0.01 billion rise in the equity value ($\partial s/\partial E = -0.0047$ per USD
billion). USD~10 million of sold five-year protection at the 500 coupon has a \omterm{def:rc:reduced-form-credit:cs01}{CS01} of
$-\text{USD}~3\,713$ per basis point at a quote of 500, so it gains USD~174\,423 per USD~1
billion of equity value: the hedge is a short of USD~1.74 million of stock
(\cref{fig:rc:structural-credit-models:equityspread}).
\end{example}

\begin{omfigure}
\begin{tikzpicture}
\begin{axis}[width=11.5cm,height=5.4cm,
  xlabel={equity value (USD billion)},ylabel={five-year model spread (bp)},
  tick label style={font=\small},label style={font=\small},
  xmin=6,xmax=16,ymin=0,ymax=500,grid=major,grid style={gray!25},
  table/col sep=comma]
\addplot[omDef,very thick] table[x=e,y=spread]{figdata/rates-credit-risk/14-structural-credit-models/equity_spread.csv};
\addplot[omThm,only marks,mark=*,mark size=2.2pt] coordinates {(10,228.7) (11.8,160.7)};
\node[omThm,font=\footnotesize,anchor=south west] at (axis cs:10,232) {before: 229 bp};
\node[omThm,font=\footnotesize,anchor=south west] at (axis cs:11.8,164) {after +18\%: 161 bp};
\end{axis}
\end{tikzpicture}

\omcaption{The \omterm{def:rc:structural-credit-models:fp}{first-passage model}'s five-year spread against the firm's equity value,
asset volatility held (barrier 70\% of face). The slope at USD~10 billion is the hedge
ratio; the curve's convexity is why an 18\% move leaves a hedged position with a loss even
if the model is right. Data: the chapter's tutorial.}\label{fig:rc:structural-credit-models:equityspread}
\end{omfigure}

\begin{remark}[What the model assumes about volatility]\label{rem:rc:structural-credit-models:vol}
Holding the asset volatility fixed makes the equity volatility fall as equity rises
(leverage falls): after an 18\% rally it is 47.1\%. If instead the equity volatility stays at
50\%, the inversion gives a higher asset volatility, 17.3\%, and the model spread
\emph{rises}, to 251 basis points. The sign of the hedge depends on which volatility the
desk believes is stable, and on equity options, if the firm has them.
\end{remark}

\section{Tutorial: from equity to spreads}\label{tut:rc:structural-credit-models:tutorial}

\begin{tutorial}
\textbf{Goal.} Back out a firm's assets from its equity, measure the \omterm{def:rc:structural-credit-models:dd}{distance to default},
build model spreads in the Merton and \omterm{def:rc:structural-credit-models:fp}{first-passage models}, and size an equity hedge.
\textbf{End state:} the numbers of
\cref{ex:rc:structural-credit-models:firm,ex:rc:structural-credit-models:fp,ex:rc:structural-credit-models:hedge}
and \cref{fig:rc:structural-credit-models:payoff,fig:rc:structural-credit-models:merton,fig:rc:structural-credit-models:fp,fig:rc:structural-credit-models:equityspread}.
\begin{tutsteps}
\item \textbf{Invert}: \texttt{invert(10, 0.5, 40, 5, 0.04)}; check equity and equity
volatility reprice.
\item \textbf{Merton}: \texttt{firm()} for the spread and both probabilities of default.
\item \textbf{First passage}: \texttt{fp\_spread} for each maturity and
\texttt{implied\_barrier()} for the 500-basis-point quote.
\item \textbf{Hedge}: \texttt{hedge()}; \texttt{fig\_rc\_structural.py} writes the charts.
\end{tutsteps}
\textbf{What to change next.} Let the barrier grow at the risk-free rate
(\texttt{growth}) and compare the curves; recompute the hedge with the equity volatility
held instead of the asset volatility.
\end{tutorial}

\section{Build: structural models}\label{bld:rc:structural-credit-models:structural}

\begin{build}
\bfield{Purpose} The firm's equity-to-credit map: model spreads for names with listed
equity, the equity hedge of credit positions, and a check on default-swap marks.
\bfield{Interface} \texttt{merton\_equity}, \texttt{merton\_debt}, \texttt{merton\_spread},
\texttt{merton\_pd}, \texttt{distance\_to\_default}, \texttt{equity\_vol}; \texttt{invert(E,
sigma\_E, D, T, r)}; \texttt{asset\_from\_equity}; \texttt{first\_passage\_survival};
\texttt{hazard\_curve\_from\_survival}; \texttt{model\_cds\_spread(q, r, maturity,
recovery)}; \texttt{FlatDisc}.
\bfield{Rules} Flat rate; one debt face; survival functions become hazard curves on the
quarterly grid and are priced by chapter~13's \texttt{firm\_cdscurve}.
\bfield{Acceptance tests} \texttt{code/firm/structural/tests/}: the inversion reprices
equity and its volatility; debt plus equity is the assets and debt is below the riskless
value; the first-passage formula reduces to the reflection principle without drift and to
one with a vanishing barrier; a barrier defaults more than a terminal test; a flat
survival gives chapter~13's par spread.
\bfield{Stretch} Barrier with uncertain level (a random recovery of the default point);
jumps in the assets; a term structure of debt; the equity-option calibration of asset
volatility.
\end{build}

\begin{omsources}
\item R.\ C.\ Merton, ``On the pricing of corporate debt: the risk structure of interest
rates'', \emph{Journal of Finance} 29(2), 1974, 449--470.
\item F.\ Black and J.\ C.\ Cox, ``Valuing corporate securities: some effects of bond
indenture provisions'', \emph{Journal of Finance} 31(2), 1976, 351--367.
\item F.\ Yu, ``How profitable is capital structure arbitrage?'', \emph{Financial Analysts
Journal} 62(5), 2006, 47--62.
\item C.\ C.\ Finger (ed.), \emph{CreditGrades Technical Document}, RiskMetrics Group, May 2002.
\item Bank for International Settlements, \emph{Quarterly Review}, June 2005, box ``Stress
testing of credit markets: the downgrade of General Motors and Ford''.
\end{omsources}

\section{Exercises}

\begin{exercise}[$\star$]\label{exo:rc:structural-credit-models:1}
Assets 50, debt face 40 due in five years, asset volatility 20\%, $r = 4\%$. Give the
Merton equity and debt values, the yield spread and the risk-neutral \omterm{def:rc:reduced-form-credit:pd}{probability of
default}.
\end{exercise}

\begin{exercise}[$\star$]\label{exo:rc:structural-credit-models:2}
Why are Merton spreads near zero at short maturities for a healthy firm?
\end{exercise}

\begin{exercise}[$\star$]\label{exo:rc:structural-credit-models:3}
Give the one-year \omterm{def:rc:structural-credit-models:dd}{distance to default} and the real-world \omterm{def:rc:reduced-form-credit:pd}{probability of default} of the
chapter's firm, with the same debt face and drift.
\end{exercise}

\begin{exercise}[$\star\star$]\label{exo:rc:structural-credit-models:4}
Without drift, a firm's assets are 1.5 times the barrier with an asset volatility of 30\%.
Give the probability that it survives two years.
\end{exercise}

\begin{exercise}[$\star\star$]\label{exo:rc:structural-credit-models:5}
Size the equity hedge of USD~20 million of sold protection on the chapter's firm.
\end{exercise}

\begin{exercise}[$\star\star$]\label{exo:rc:structural-credit-models:6}
Why is equity much more volatile than the assets, and why does the ratio change when the
equity price moves?
\end{exercise}

\begin{exercise}[$\star\star\star$]\label{exo:rc:structural-credit-models:7}
\emph{Coding.} Find the barrier that fits a five-year quote of 500 basis points, and give
the model's ten-year spread with that barrier.
\end{exercise}

\begin{exercise}[$\star\star\star$]\label{exo:rc:structural-credit-models:8}
\emph{Find the flaw.} ``The \omterm{def:rc:structural-credit-models:merton}{Merton model} gives this firm 130 basis points and the market
quotes 500. The default swap is 370 basis points too wide: sell protection.''
\end{exercise}

\section{Problem: The Capital-Structure Trade}

\begin{problem}[{Weekend problem --- May 2005}]\label{pb:rc:structural-credit-models:1}
A fund sees the chapter's firm quoted at 500 basis points at five years while its
\omterm{def:rc:structural-credit-models:fp}{first-passage model} (barrier 70\% of the face, fitted on the firm's history) gives 229. It
sells USD~10 million of five-year protection at the 500 coupon and shorts equity by the
model's hedge ratio. The next day the equity rises 18\%; the day after, the firm is
downgraded and the market spread widens by 200 basis points, to 700.

\medskip\noindent\textbf{Part I --- The trade.}
\begin{enumerate}
\item Give the model spread and the market spread, and the signal.
\item Give the hedge ratio and the equity short.
\item What does the fund earn while nothing happens?
\item What view is the fund taking, stated without the model?
\item Which risks does the hedge leave open?
\end{enumerate}

\medskip\noindent\textbf{Part II --- The two days.}
\begin{enumerate}[resume]
\item Give the model spread after the 18\% rally.
\item Give the market spread the model predicted after the rally, applying the model's
change to the quote.
\item Give the equity leg's P\&L.
\item Give the default swap's P\&L had the market followed the model, and the hedged P\&L.
\item Give the default swap's P\&L at 700 and the total.
\end{enumerate}

\medskip\noindent\textbf{Part III --- What the model missed.}
\begin{enumerate}[resume]
\item Give the hedge ratio after the rally.
\item Recompute the model spread after the rally with the equity volatility held at 50\%.
\item Why can an equity rally and a spread widening happen together?
\item At what market spread would the trade have broken even on the two days?
\item Which instrument would have hedged the downgrade?
\end{enumerate}

\medskip\noindent\textbf{Part IV --- Judgement.}
\begin{enumerate}[resume]
\item Why did many funds lose on similar trades in the same week?
\item How would you size such trades?
\item What would make you close the trade after the two days, and what would make you
keep it?
\item State the \emph{named result}: the two legs' P\&L and the total on the two days.
\item In one sentence: what does a \omterm{def:rc:structural-credit-models:structural}{structural model} add to a credit desk?
\end{enumerate}
\end{problem}

\section{Interview questions}

\begin{interviewq}[$\star$ \iqroles{researcher, trader}]\label{iq:rc:structural-credit-models:1}
Why is equity a call option on the firm? What is the debt?
\end{interviewq}

\begin{interviewq}[$\star\star$ \iqroles{researcher, developer}]\label{iq:rc:structural-credit-models:2}
How do you estimate a firm's asset value and asset volatility?
\end{interviewq}

\begin{interviewq}[$\star\star$ \iqroles{researcher}]\label{iq:rc:structural-credit-models:3}
Why do \omterm{def:rc:structural-credit-models:structural}{structural models} underprice short-dated credit, and how do you fix it?
\end{interviewq}

\begin{interviewq}[$\star\star$ \iqroles{risk, bank}]\label{iq:rc:structural-credit-models:4}
What is a \omterm{def:rc:structural-credit-models:dd}{distance to default}, and how is it turned into a \omterm{def:rc:reduced-form-credit:pd}{probability of default} in
practice?
\end{interviewq}

\begin{interviewq}[$\star\star\star$ \iqroles{trader}]\label{iq:rc:structural-credit-models:5}
Describe a \omterm{def:rc:structural-credit-models:csa}{capital-structure arbitrage} trade and its hedge ratio. How can it lose on
both legs?
\end{interviewq}

\begin{interviewq}[$\star\star\star$ \iqroles{researcher, trader}]\label{iq:rc:structural-credit-models:6}
Compare the reduced-form and structural approaches. When would you use each?
\end{interviewq}
